% !PRINTED 12
\weberpage{12}

Wir erwähnen noch folgenden Satz, dessen Richtigkeit sich unmittelbar aus
dem Isomorphismus transformirter Gruppen ergiebt:

\gesperrt{Ist $Q$ ein Theiler von $P$, $R$ ein Theiler von $Q$, so ist, wenn $Q'$
und $R'$ aus $Q$ und $R$ durch dasselbe Element von $P$ transformirt
sind, auch $R'$ ein Theiler von $Q'$, und wenn $R$ Normaltheiler von $Q$
ist, so ist auch $R'$ Normaltheiler von $Q'$.}

\begin{center}
\S.\,4.

\gesperrt{Composition der Theile.}
\end{center}

Sind $A$ und $B$ irgend zwei Reihen von Elementen aus einer Gruppe $P$
(Gruppen oder nicht), so wollen wir unter dem symbolischen Producte $AB$
den Inbegriff aller Elemente verstehen, die man erhält, wenn man je ein
Element $a$ von $A$ mit je einem Element $b$ von $B$ nach der in $P$
geltenden Vorschrift zu $ab$ componirt. Diese Art der Zusammensetzung von
$A$ und $B$ zu $AB$ wollen wir die \gesperrt{Composition der Theile} (von $P$)
nennen. Wir unterscheiden hier zwischen einem \gesperrt{Theil} und einem \gesperrt{Theiler} von
$P$, so dass ein Theiler immer eine Gruppe sein soll, was bei einem Theil
nicht nothwendig ist.

Man kann ebenso drei und mehr Theile $A, B, C \ldots$ von $P$ componiren,
und es gilt bei der Composition der Theile das associative Gesetz, wie
unmittelbar daraus folgt, dass dieses Gesetz in $P$ gilt. Danach ist die
Bedeutung eines Compositums $ABC \ldots$ eindeutig bestimmt.

In der Zusammensetzung $AB$ kann einer der Theile $A$, $B$ auch aus einem
einzigen Elemente bestehen, und dann stimmt die Bezeichnung $AB$ mit der
oben für die Nebengruppen gebrauchten Bezeichnung überein.

Bei der Composition der Theile gelten folgende Sätze:

\noindent 1.\ \gesperrt{Die nothwendige und hinreichende Bedingung
dafür, dass ein Theil $A$ von $P$ eine Gruppe ist, besteht in der
Gleichung}
\weq{1}{AA = A.}
Denn diese Gleichung besagt nichts Anderes als dass, wenn $a$ und $a'$ in
$A$ enthalten sind, auch $aa'$ zu $A$ gehört, also dass $A$ eine Gruppe
ist.
